\makeabstract
{
Let’s Get Metal: An Analysis of Gas-Phase Metallicity for Galaxies within the DIVE Survey
}
{
Meadow Perez (1), Mia de los Reyes(1), Leonardo Barba (2,1)
}
{
\textsuperscript{1}Department of Physics and Astronomy, Amherst College, Amherst, MA, USA\\
\textsuperscript{2}Department of Astronomy, University of Massachusetts Amherst, Amherst, MA, USA
}
{
This analysis looks at the metal abundances of dwarf galaxies living in void environments. Using Integral Field Unit (IFU) data, this analysis aims to develop a new method of measuring metal in a galaxies’ gas and utilize different calibration methods in order to discover possible trends in a Gas-Phase Metallicity (GPM) plot.
}
{
This analysis used the data obtained from de los Reyes et al. (2023) through the KCWI spectrograph. Spectral data were gathered through a method of weighted integration which allows the entire galaxy to be included. Using the documented and appropriate GPM calibration methods outlined in Kewley \& Ellison (2008), the appropriated spectral lines indicating certain elements within the gas were able to be isolated. The lines were then fitted with a gaussian distribution and integrated in order to measure the total flux. With this flux, their GPM’s were calculated and plotted. Due to a lack of data (5 control galaxies, 16 void galaxies), no statistical analysis was performed.
}
{
The results of this analysis show that IFU data have promising uses on the further study of intergalactic objects, especially galaxies. This is one of the first times full-body representation has been considered in looking at galaxy spectra and can provide more accurate analysis. As can be seen in the final GPM plots, across all calibrations, none of the DIVE galaxies seem to indicate a trend. While this does not rule that a trend won’t emerge upon the collection of more data, it does validate the accuracy of spectral integration. Further research would include the analysis of dwarf galaxies to find statistical significance of any trends, and the analysis of large galaxies to confirm if spectral integration in IFU data matches known GPM trends.
}
{
In conclusion, this analysis provides strong evidence for the accuracy and use of IFU data in spectral analysis of objects. The similarities between the GPM plots suggests a consistent measurement of metals with galaxies. This validates the use and need for more images under the IFU format.
}

\newpage
